% Figures F1-F7. F1/F4/F5 are deliberately textual (they reproduce
% real compiler output / real IR states verbatim -- fact-ledgered in
% paper/figures/); F2/F3 are TikZ from the ledgered sources; F6 is a
% placeholder pending the capture protocol; F7 is the generated,
% self-checking SVG (tools/plot_overhead.py).

\begin{figure*}[t]
\begin{lstlisting}
ERROR TYPE:                          |  HOW TO FIX:
Malformed Expression                 |  Insert an operator between x and 2
                                     |  or remove the extra value
LOCATION:                            |
line 3, column 14                    |  TRACE:
    3 |     return x 2;              |  Parser
      |              ^               |  -> Parser::parse()
                                     |  -> Parser::parseFunctionDecl()
ROOT CAUSE:                          |       [starting at line 1]
Two expressions appeared             |  -> Parser::parseBlock()
consecutively with no operator       |  -> Parser::parseStatement()
between them.                        |       [at 'return' (line 3)]
                                     |  -> Parser::parseReturnStmt()
WHAT HAPPENED:                       |       [at 'return' (line 3)]
The parser successfully read:        |  -> DiagnosticEngine::
    return x                         |       malformedExpression()
After completing that expression,    |       [diagnostic created]
it encountered another value:        |
    2                                |  INVALID:   x 2
C++ requires an operator             |  VALID:     x + 2
between values.                      |             x * 2
                                     |             x = 2
\end{lstlisting}
\caption{One diagnostic, verbatim (two-column arrangement of the
CLI's sequential output). Every field is data on one struct
(\S\ref{sec:architecture}): the trace is recorded from the parser's
actual descent with runtime details in brackets, the examples are
synthesized from the offending tokens, and the same value renders
identically in the visualizer (Fig.~\ref{fig:f6}).}
\label{fig:f1}
\end{figure*}

\begin{figure}[t]
\centering
\begin{tikzpicture}[node distance=3.5mm and 5mm, >={Stealth},
  st/.style={draw, rounded corners=1pt, font=\scriptsize, inner sep=2.5pt},
  dg/.style={draw, rounded corners=1pt, font=\scriptsize, inner sep=2.5pt, fill=black!8},
  e/.style={->, font=\tiny}]
\node[st] (lx) {Lexer};
\node[st, right=of lx] (ps) {Parser};
\node[st, right=of ps] (sa) {Semantic};
\node[st, right=of sa] (ir) {IR};
\node[st, right=of ir] (op) {Optimizer};
\node[st, right=of op] (cg) {Codegen};
\node[dg, below=7mm of ps, xshift=8mm] (dc) {DiagnosticCollector};
\node[st, below=3.5mm of dc, xshift=-14mm] (r1) {\tiny full text};
\node[st, below=3.5mm of dc] (r2) {\tiny compact};
\node[st, below=3.5mm of dc, xshift=13mm] (r3) {\tiny JSON};
\draw[e] (lx) -- (ps); \draw[e] (ps) -- (sa); \draw[e] (sa) -- (ir);
\draw[e] (ir) -- (op); \draw[e] (op) -- (cg);
\draw[e, dashed] (lx) |- (dc);
\draw[e, dashed] (ps) -- (dc);
\draw[e, dashed] (sa) |- (dc);
\draw[e] (dc) -- (r1); \draw[e] (dc) -- (r2); \draw[e] (dc) -- (r3);
\draw[e, dotted] (ir) to[bend left=18] node[above, font=\tiny]{span+note} (cg);
\end{tikzpicture}
\caption{The pipeline with its diagnostics spine. Every front-end
stage returns output \emph{and} diagnostics (D1); IR, optimizer, and
codegen emit none by design (D2) --- but carry provenance (dotted).
Three renderers consume one struct (D6).}
\label{fig:f2}
\end{figure}

\begin{figure}[t]
\centering
\begin{tikzpicture}[node distance=3mm and 4mm, >={Stealth},
  r/.style={draw, rounded corners=1pt, font=\scriptsize, align=center, inner sep=2.5pt},
  e/.style={->, font=\tiny}]
\node[r] (k) {Kind taxonomy\\ \tiny 17 kinds, 5 reserved};
\node[r, right=of k] (kb) {Knowledge base\\ \tiny prose per root cause};
\node[r, right=of kb] (tr) {TraceRecorder\\ \tiny per stage, RAII};
\node[r, below=6mm of kb] (en) {DiagnosticEngine\\ \tiny 16 factory methods};
\node[r, below=5mm of en] (d) {Diagnostic value};
\node[r, below=5mm of d] (c) {Collector $\to$ text / compact / JSON};
\draw[e] (k) -- (en); \draw[e] (kb) -- (en);
\draw[e] (tr) -- node[right, font=\tiny]{snapshot} (en);
\draw[e] (en) -- (d); \draw[e] (d) -- (c);
\end{tikzpicture}
\caption{The five-role diagnostic core (\S\ref{sec:architecture});
the fifth role, the \texttt{StageOutput} transport, is Fig.~\ref{fig:f2}'s
spine. The knowledge-base edge also carries the curated fallback
chains, unreachable in the pipeline (D8).}
\label{fig:f3}
\end{figure}

% F4 is folded into F1's trace panel for space at first assembly; a
% dedicated sequence diagram is cut unless a reviewer asks. The
% ledgered source remains in paper/figures/F2-F5-diagram-sources.md.

\begin{figure}[t]
\begin{lstlisting}
source (line 2):   return (2 + 3) * 4;

before        iteration 1          iteration 2         final
[2] t0 = 2+3  CF: t0 = 5           (t0 gone)
                 ;folded from
                  't0 = 2 + 3'
[2] t1 = t0*4 CP: t1 = 5 * 4       CF: t1 = 20         (t1 gone -- its
                 ;t0->5 (CopyProp)    ;t0->5; folded     fold AND prop
              DCE: t0 removed          from 't1 = 5*4'   notes die)
                  (fold note dies) CP: return 20
[2] return t1     return t1           ;t1->20 (CopyProp)
                                   DCE: t1 removed     [2] return 20
                                                          ;t1->20 (CopyProp)
\end{lstlisting}
\caption{Provenance through the fixed point for \texttt{(2+3)*4}.
The source line survives every rewrite; transformation notes
accumulate --- and die twice with the instructions DCE removes. The
single surviving instruction reports only the propagation step:
per-instruction provenance is honest but not transitive
(\S\ref{sec:limitations}, item 3).}
\label{fig:f5}
\end{figure}

\begin{figure}[t]
\centering
\fbox{\parbox{0.9\linewidth}{\centering\scriptsize\vspace{8mm}
F6: visualizer screenshot --- capture per
\texttt{paper/figures/F6-F7-notes.md} (full-config build; same
program family as Fig.~\ref{fig:f1}).\vspace{8mm}}}
\caption{The visualizer rendering the same diagnostic value as
Fig.~\ref{fig:f1}: identical sections, identical recorded trace,
because both surfaces are renderers over one struct (D6). Right: IR
panel with per-instruction source-line gutters and optimizer notes.}
\label{fig:f6}
\end{figure}

\begin{figure*}[t]
\centering
\includesvg[width=0.95\textwidth]{figures/F7-overhead}
\caption{Always-on overhead of the explanation architecture versus
the ablated baseline, before and after the lazy-formatting repair
(D14), across program sizes (log scale). Whiskers: 95\% bootstrap
CIs. Left: eager detail strings cost up to $\sim$175\% (trace) and
$\sim$160\% (full). Right, same scale: recording collapses to CIs
spanning zero; the remaining 12--17\% is provenance
(\S\ref{sec:eval}). Generated by a script that hard-asserts every
plotted median against the paper's tables.}
\label{fig:f7}
\end{figure*}
